%%%%%%%%%%%%%%%%%%%%%%%%%%%%%%%%%%%%%%%%%%%%%%%%%%%%%%%%%%%%%%%%%%%%%%%%%%%%%

\section*{Remerciements}

En premier lieu, je souhaiterais exprimer toute ma reconnaissance à François-Gildas
Tual pour la qualité de ses cours, pour sa bienveillance envers les élèves ainsi
que pour ses commentaires.

J'adresse ma gratitude à Olga Onegina pour m’avoir expliqué où trouver de
l'information en langue russe et je remercie Olivier Morlotti pour sa relecture
de mon mémoire.

%%%%%%%%%%%%%%%%%%%%%%%%%%%%%%%%%%%%%%%%%%%%%%%%%%%%%%%%%%%%%%%%%%%%%%%%%%%%%

\vspace{4.0cm}

%%%%%%%%%%%%%%%%%%%%%%%%%%%%%%%%%%%%%%%%%%%%%%%%%%%%%%%%%%%%%%%%%%%%%%%%%%%%%

\section*{Contenu multimédia associé}

Dans ce mémoire, chaque \emph{Quick Response code} (QR code) indique la
présence d'un contenu multimédia accessible à l'aide d'un smartphone.

\begin{figure}[!ht]
  \begin{bigcenter}
    \includegraphics[width=3cm, keepaspectratio]{qr-cover.pdf}
  \end{bigcenter}
  \caption{Exemple de QR code.}
\end{figure}

%%%%%%%%%%%%%%%%%%%%%%%%%%%%%%%%%%%%%%%%%%%%%%%%%%%%%%%%%%%%%%%%%%%%%%%%%%%%%
